\markboth{AGENT INTERACTION}{AGENT INTERACTION}
\section{Agent Interaction and Communication}\label{sec:agents}

This section details how the seven agents coordinate to answer user queries: communication protocols, delegation patterns, prompt architecture, the recording/clipping strategy for large tool results, and observability infrastructure.

\subsection{Communication Protocols}\label{sec:agents:protocols}

Training Copilot employs two JSON-RPC\,2.0-based protocols at different layers. \textbf{MCP} governs all data access: each specialist calls its servers via a scoped \texttt{ToolHost} over Streamable HTTP, enforced at process start: \texttt{scoped\_host(["garmin"])} builds a \texttt{ToolHost} containing only the Garmin URL, so the specialist's discovery call returns only Garmin tools and cannot even see the Strava or Weather tools, regardless of its prompt. This is a stronger guarantee than prompt-based access control. \textbf{A2A} governs inter-agent communication: the orchestrator calls each specialist via \texttt{ask\_<name>} tools, each a streaming A2A request through the \texttt{a2a-sdk}. \texttt{call\_agent()} resolves the specialist's Agent Card at \texttt{/.well-known/agent-card.json} and sends the question; the response is a stream of typed events: status updates shown in the UI as progress (``Consulting recovery agent\ldots''), \texttt{TextPart}s (the answer), and \texttt{DataPart}s (the full tool-call records as JSON). This separation is central: the user-facing answer stays concise for the orchestrator's context, while the data artifact carries complete MCP results (GPS streams, intraday timelines) for the UI to render maps and charts. Figure~\ref{fig:interaction_flow} illustrates the flow for a representative multi-domain query.

\begin{figure}[ht]
\centering
\resizebox{\textwidth}{!}{%
\begin{tikzpicture}[
    node distance=0.3cm,
    every node/.style={font=\small},
    actor/.style={draw, thick, minimum width=2.0cm, minimum height=0.6cm, rounded corners=2pt, align=center},
    msg/.style={-{Stealth[length=4pt]}, thick},
    rmsg/.style={-{Stealth[length=4pt]}, thick, dashed},
    timelabel/.style={font=\scriptsize\itshape, text=tcgrey!50},
]

    \node[actor, draw=tcblue, fill=tcblue!10] (user) at (0, 0) {User};
    \node[actor, draw=tcgreen, fill=tcgreen!8] (orch) at (3.2, 0) {Orchestrator};
    \node[actor, draw=tcyellow!80!black, fill=tcyellow!10] (rec) at (6.4, 0) {Recovery};
    \node[actor, draw=tccyan, fill=tccyan!8] (ctx) at (9.6, 0) {Context};

    % Lifelines
    \draw[thick, tclightgrey] (user.south) -- ++(0, -8.5);
    \draw[thick, tclightgrey] (orch.south) -- ++(0, -8.5);
    \draw[thick, tclightgrey] (rec.south) -- ++(0, -8.5);
    \draw[thick, tclightgrey] (ctx.south) -- ++(0, -8.5);

    % 1. User → Orchestrator
    \draw[msg, color=tcblue] (0, -1.0) -- node[above, font=\scriptsize] {``Should I train tomorrow?''} (3.2, -1.0);
    \node[timelabel] at (-1.3, -1.0) {t\textsubscript{0}};

    % 2. Orchestrator → Recovery + Context (parallel)
    \draw[msg, color=tcyellow!80!black] (3.2, -2.0) -- node[above, font=\scriptsize] {\texttt{ask\_recovery}} (6.4, -2.0);
    \draw[msg, color=tccyan] (3.2, -2.3) -- node[below, font=\scriptsize] {\texttt{ask\_context}} (9.6, -2.3);
    \node[timelabel] at (-1.3, -2.15) {t\textsubscript{1} (parallel)};

    % 3. Recovery → Garmin MCP (internal)
    \node[font=\scriptsize, text=tcgrey, align=left] at (7.2, -3.2) {\texttt{garmin\_\_sleep}\\[1pt]\texttt{garmin\_\_hrv}\\[1pt]\texttt{garmin\_\_body\_battery}};

    % 4. Context → Weather + Calendar MCP (internal)
    \node[font=\scriptsize, text=tcgrey, align=left] at (10.4, -3.7) {\texttt{weather\_\_forecast}\\[1pt]\texttt{calendar\_\_events}};

    % 5. Recovery → Orchestrator (response)
    \draw[rmsg, color=tcyellow!80!black] (6.4, -5.0) -- node[above, font=\scriptsize] {answer + DataPart} (3.2, -5.0);

    % 6. Context → Orchestrator (response)
    \draw[rmsg, color=tccyan] (9.6, -5.5) -- node[above, font=\scriptsize] {answer + DataPart} (3.2, -5.5);
    \node[timelabel] at (-1.3, -5.25) {t\textsubscript{2}};

    % 7. Orchestrator synthesises
    \node[font=\scriptsize, text=tcgreen, align=center] at (3.2, -6.5) {Synthesise\\+ build trace};

    % 8. Orchestrator → User
    \draw[rmsg, color=tcblue] (3.2, -7.5) -- node[above, font=\scriptsize, align=center] {``HRV 12\% below baseline\ldots\\14--17h rain-free. Light session.''} (0, -7.5);
    \node[timelabel] at (-1.3, -7.5) {t\textsubscript{3}};

\end{tikzpicture}%
}% end resizebox
\caption[Sequence diagram for a multi-domain query]{Sequence diagram for the query ``Should I train tomorrow?''. The orchestrator delegates to recovery and context specialists in parallel; each calls its scoped MCP servers, then returns an answer with a data artifact.}\label{fig:interaction_flow}
\end{figure}


\subsection{Delegation Patterns}\label{sec:agents:delegation}

The orchestrator employs four delegation patterns. \textit{Single-specialist delegation} routes a query cleanly to one domain (e.g.\ ``What's my VO\textsubscript{2}max?'' $\to$ load specialist). \textit{Parallel multi-specialist delegation}, the most common pattern, emits multiple \texttt{ask\_*} calls in a single LLM step, reducing latency to the slowest specialist rather than their sum. \textit{Sequential delegation} occurs when a later specialist depends on an earlier one's output, which the LLM handles naturally across multiple reasoning steps. The fourth, \textit{schedule-triggered delegation}, follows the same three shapes above but originates from a background scheduler rather than a live chat message: on a fixed cadence, scheduled check-ins (a daily one, plus nudges around calendar events and the coach's own \texttt{schedule\_followup} entries, Section~\ref{sec:architecture:supporting}) are composed by running their stored note through the same specialist layer, and the resulting summary is pushed proactively to the in-app Coach chat and the Telegram bridge instead of being returned to an open chat request.

\subsection{Prompt Architecture}\label{sec:agents:prompts}

The prompt system, defined entirely in \texttt{agents/prompts.py}, composes three layers, following the task-specific instruction design that prompt-engineering research~\citep{sahoo_prompt_2024} identifies as important for reliable behaviour without changing model weights. A shared \textbf{base prompt} gives every agent the same identity (a supportive, concise training buddy who holds the athlete accountable without being either a clinical assistant or a drill sergeant), injects the current date (computed at call time), sets the home location for geocoding defaults, and establishes five rules: use tools for any data question and never fabricate; execute tool calls immediately (``the question IS the permission''); call independent tools in parallel; compute absolute dates before passing them to tools; and synthesise data into insight rather than dumping raw JSON. \textbf{Specialist prompts} add a domain block that lists the available tools with explicit \textit{when}-to-use guidance. For example, the load specialist maps ``weekly volume'' to \texttt{strava\_\_get\_training\_trends} and sets Strava as the primary activity source with Garmin as fallback, which avoids contradictory answers from duplicated data; the recovery specialist adds rules like ``interpret HRV relative to personal baseline; flag overtraining when HRV is suppressed, Body Battery low, and sleep poor''; the coach specialist is told to let the athlete server do all arithmetic (zones, plan volumes, feasibility) and never to invent such numbers itself. The \textbf{orchestrator prompt} defines the routing policy (always delegate, never answer from parametric knowledge, pick the minimal specialist set, and preserve the fitness specialist's citations verbatim) and lists only specialist domains, never individual tools.

\subsection{Recording, Clipping, and Trace Assembly}\label{sec:agents:recording}

Tool results can be large (a GPS stream from a two-hour ride holds thousands of points, and a 14-day wellness trend is a multi-kilobyte JSON array), while the LLM's context window is finite and every token costs~\citep{qu_tool_2024}. Training Copilot handles this with a \textit{record-full, clip-for-model} pattern (Figure~\ref{fig:recording_clipping}), implemented in the \texttt{build\_tools} wrapper (\texttt{core/mcp\_langchain.py}): each tool call appends its full JSON result to a shared \texttt{recorder} list with the tool name, arguments, duration, and any error; the \texttt{clip()} function then replaces arrays under specific keys (\texttt{points}, \texttt{waypoints}, \texttt{timeline}, \ldots) with a count placeholder (e.g.\ \texttt{[847 items]}), truncates long strings, and caps the total at 6{,}000 characters. The clipped string is what the LLM sees. The UI renders maps and charts from the \textit{full} trace, not the model's response, decoupling data fidelity from context-window constraints.

\begin{figure}[ht]
\centering
\resizebox{\textwidth}{!}{%
\begin{tikzpicture}[
    node distance=0.4cm and 0.5cm,
    every node/.style={font=\small},
    block/.style={draw, thick, rounded corners=3pt, minimum height=0.8cm, minimum width=2.2cm, align=center},
    data/.style={draw, thick, rounded corners=1pt, minimum height=0.55cm, align=center, font=\scriptsize},
    arr/.style={-{Stealth[length=5pt]}, thick},
]
    \node[block, draw=tcgrey!60, fill=tcgrey!5] (mcp) {MCP Server};
    \node[block, draw=tcgreen, fill=tcgreen!6, right=1.8cm of mcp] (bridge) {MCP-LangChain\\Bridge};

    \node[data, draw=tcgreen, fill=tcgreen!8, above right=0.7cm and 1.5cm of bridge] (full) {\textbf{Full result}\\(847 GPS points)};
    \node[data, draw=tcyellow!80!black, fill=tcyellow!8, below right=0.7cm and 1.5cm of bridge] (clipped) {\textbf{Clipped copy}\\(\texttt{[847 items]})};

    \node[block, draw=tcgreen!80, fill=tcgreen!6, right=1.2cm of full] (trace) {Trace /\\DataPart};
    \node[block, draw=tcyellow!80!black, fill=tcyellow!6, right=1.2cm of clipped] (llm) {LLM\\Context};

    \node[block, draw=tcblue, fill=tcblue!8, right=0.8cm of trace] (ui) {UI Renderer\\(maps, charts)};
    \node[block, draw=tcred!80, fill=tcred!6, right=0.8cm of llm] (answer) {Agent\\Reasoning};

    \draw[arr, color=tcgrey] (mcp) -- node[above, font=\scriptsize]{JSON} (bridge);
    \draw[arr, color=tcgreen] (bridge) -- node[left, font=\scriptsize, rotate=30]{record} (full);
    \draw[arr, color=tcyellow!80!black] (bridge) -- node[left, font=\scriptsize, rotate=-30]{clip} (clipped);
    \draw[arr, color=tcgreen] (full) -- (trace);
    \draw[arr, color=tcyellow!80!black] (clipped) -- (llm);
    \draw[arr, color=tcblue] (trace) -- (ui);
    \draw[arr, color=tcred!70] (llm) -- (answer);

\end{tikzpicture}%
}% end resizebox
\caption[Recording and clipping pattern]{The recording and clipping pattern. Full tool results are recorded into the trace artifact for UI rendering, while a clipped copy (large arrays replaced with placeholders, capped at 6\,000 characters) is returned to the LLM's context for reasoning. This decoupling preserves data fidelity for rendering without wasting LLM context tokens.}\label{fig:recording_clipping}
\end{figure}


When a specialist finishes, its accumulated \texttt{recorder} list attaches to the A2A response as a \texttt{DataPart} artifact. The orchestrator collects these artifacts from all consulted specialists and passes them to \texttt{build\_trace()}, which assembles a structured \textit{trace}: run ID, timestamp, original input, every tool call with arguments, full results, duration and error status, agents consulted, extracted route data (the first successful result from a route-producing tool), chart hints extracted from \texttt{<!--charts: ...-->} tags the model appends, and the final answer. This trace serves transparency (the chat UI's collapsible debug panel, Figure~\ref{fig:screenshot_chat} in Appendix~\ref{sec:appendix:screenshots}), rendering (route maps and inline charts from full data), and evaluation (the \texttt{grounded\_in\_real\_data} scorer of Section~\ref{sec:evaluation:scorers} reads tool spans directly from it).

\subsection{Observability and Graceful Degradation}\label{sec:agents:observability}

All agent runs are traced to an MLflow tracking server~\citep{mlflow_multiturn_2025}. Each process calls \texttt{setup\_tracing(name)} at startup, enabling LangChain autologging and wrapping each \texttt{ainvoke} in a labelled root span, producing a hierarchical span tree per run filterable by service, with LLM and tool calls as child spans; tracing is strictly best-effort, so a missing \texttt{mlflow} or unreachable server logs once and is ignored. Execution traces answer \textit{what} the system did, but not whether the user considered the answer wrong, unhelpful, or unsafe; a one-click \textit{report button} closes this loop. Pressing it captures a diagnostic bundle (recent logs, MLflow trace references, and the reporting user's own state: profile, goals, schedules, and recent chats) and notifies the maintainers, so a flagged response can be cross-referenced against its full agent trace during later analysis. Turning this free-text feedback into a labelled good/bad validation set for the LLM judges (Section~\ref{sec:evaluation}) is left as future work. Graceful degradation is systematic at every layer: unreachable MCP servers are skipped during discovery, unavailable specialists return descriptive errors rather than failing the task, tool errors come back as structured JSON the agent can reason about, and MLflow tracing never blocks the critical path. So Training Copilot still gives useful answers when some data sources or agents are offline, which matters for a system that depends on several external APIs.
